\uuid{riyB}
\exo7id{1941}
\titre{exo7 1941}
\auteur{gineste}
\organisation{exo7}
\datecreate{2001-11-01}
\isIndication{false}
\isCorrection{false}
\chapitre{Série numérique}
\sousChapitre{Série à  termes positifs}
\cpge{mpsi-mp2i.integration.c, pcsi.integration.b, ptsi.integration.b, mpsi-mp2i.procedes-sommatoires-discrets.a, pcsi.series-numeriques.a, ptsi.series-numeriques.a, mpsi-mp2i.nombres-reels-suites.d, pcsi.nombres-reels-suites.c, ptsi.nombres-reels-suites.c, bcpst1.integration-segment.a, bcpst1.integration-segment.b, mp-mpi.series-numeriques-vectorielles.a, pt.series-numeriques.b, mp-mpi.series-numeriques-vectorielles.b, pc.series-numeriques, psi.series-numeriques.a, pt.series-numeriques.a, bcpst2.series-reelles, bcpst2.integrales-generalisees}
\module{Analyse}
\niveau{L1}
\difficulte{}

\contenu{
\question{
Le but de cet exercice est de montrer la convergence de l'intégrale généralisée suivante $ \displaystyle \int_0^{\infty} \frac{dx}{1 + x^4 \sin^2 x}.$

Pour cela, on considère la série de terme général $$u_n=\int_{n\pi}^{(n+1)\pi} \frac{dx}{1 + x^4 \sin^2 x} .$$
Par un changement de variable, transformer $u_n$ en
$$ u_n =\int_{0}^{\pi} \frac{dx}{1 + (n\pi + x)^4 \sin^2 x} $$
Encadrer ensuite $u_n$ par les termes de la suite $v_n$ o\`u
$$ v_n =\int_{0}^{\pi} \frac{dx}{1 + (n\pi)^4 \sin^2 x} $$
Calculer explicitement l'intégrale $v_n$ et en déduire un équivalent de $u_n$. Conclure.
}
}
